% Options for packages loaded elsewhere
\PassOptionsToPackage{unicode}{hyperref}
\PassOptionsToPackage{hyphens}{url}
\documentclass[
]{article}
\usepackage{xcolor}
\usepackage[margin=1in]{geometry}
\usepackage{amsmath,amssymb}
\setcounter{secnumdepth}{-\maxdimen} % remove section numbering
\usepackage{iftex}
\ifPDFTeX
  \usepackage[T1]{fontenc}
  \usepackage[utf8]{inputenc}
  \usepackage{textcomp} % provide euro and other symbols
\else % if luatex or xetex
  \usepackage{unicode-math} % this also loads fontspec
  \defaultfontfeatures{Scale=MatchLowercase}
  \defaultfontfeatures[\rmfamily]{Ligatures=TeX,Scale=1}
\fi
\usepackage{lmodern}
\ifPDFTeX\else
  % xetex/luatex font selection
\fi
% Use upquote if available, for straight quotes in verbatim environments
\IfFileExists{upquote.sty}{\usepackage{upquote}}{}
\IfFileExists{microtype.sty}{% use microtype if available
  \usepackage[]{microtype}
  \UseMicrotypeSet[protrusion]{basicmath} % disable protrusion for tt fonts
}{}
\makeatletter
\@ifundefined{KOMAClassName}{% if non-KOMA class
  \IfFileExists{parskip.sty}{%
    \usepackage{parskip}
  }{% else
    \setlength{\parindent}{0pt}
    \setlength{\parskip}{6pt plus 2pt minus 1pt}}
}{% if KOMA class
  \KOMAoptions{parskip=half}}
\makeatother
\usepackage{color}
\usepackage{fancyvrb}
\newcommand{\VerbBar}{|}
\newcommand{\VERB}{\Verb[commandchars=\\\{\}]}
\DefineVerbatimEnvironment{Highlighting}{Verbatim}{commandchars=\\\{\}}
% Add ',fontsize=\small' for more characters per line
\usepackage{framed}
\definecolor{shadecolor}{RGB}{248,248,248}
\newenvironment{Shaded}{\begin{snugshade}}{\end{snugshade}}
\newcommand{\AlertTok}[1]{\textcolor[rgb]{0.94,0.16,0.16}{#1}}
\newcommand{\AnnotationTok}[1]{\textcolor[rgb]{0.56,0.35,0.01}{\textbf{\textit{#1}}}}
\newcommand{\AttributeTok}[1]{\textcolor[rgb]{0.13,0.29,0.53}{#1}}
\newcommand{\BaseNTok}[1]{\textcolor[rgb]{0.00,0.00,0.81}{#1}}
\newcommand{\BuiltInTok}[1]{#1}
\newcommand{\CharTok}[1]{\textcolor[rgb]{0.31,0.60,0.02}{#1}}
\newcommand{\CommentTok}[1]{\textcolor[rgb]{0.56,0.35,0.01}{\textit{#1}}}
\newcommand{\CommentVarTok}[1]{\textcolor[rgb]{0.56,0.35,0.01}{\textbf{\textit{#1}}}}
\newcommand{\ConstantTok}[1]{\textcolor[rgb]{0.56,0.35,0.01}{#1}}
\newcommand{\ControlFlowTok}[1]{\textcolor[rgb]{0.13,0.29,0.53}{\textbf{#1}}}
\newcommand{\DataTypeTok}[1]{\textcolor[rgb]{0.13,0.29,0.53}{#1}}
\newcommand{\DecValTok}[1]{\textcolor[rgb]{0.00,0.00,0.81}{#1}}
\newcommand{\DocumentationTok}[1]{\textcolor[rgb]{0.56,0.35,0.01}{\textbf{\textit{#1}}}}
\newcommand{\ErrorTok}[1]{\textcolor[rgb]{0.64,0.00,0.00}{\textbf{#1}}}
\newcommand{\ExtensionTok}[1]{#1}
\newcommand{\FloatTok}[1]{\textcolor[rgb]{0.00,0.00,0.81}{#1}}
\newcommand{\FunctionTok}[1]{\textcolor[rgb]{0.13,0.29,0.53}{\textbf{#1}}}
\newcommand{\ImportTok}[1]{#1}
\newcommand{\InformationTok}[1]{\textcolor[rgb]{0.56,0.35,0.01}{\textbf{\textit{#1}}}}
\newcommand{\KeywordTok}[1]{\textcolor[rgb]{0.13,0.29,0.53}{\textbf{#1}}}
\newcommand{\NormalTok}[1]{#1}
\newcommand{\OperatorTok}[1]{\textcolor[rgb]{0.81,0.36,0.00}{\textbf{#1}}}
\newcommand{\OtherTok}[1]{\textcolor[rgb]{0.56,0.35,0.01}{#1}}
\newcommand{\PreprocessorTok}[1]{\textcolor[rgb]{0.56,0.35,0.01}{\textit{#1}}}
\newcommand{\RegionMarkerTok}[1]{#1}
\newcommand{\SpecialCharTok}[1]{\textcolor[rgb]{0.81,0.36,0.00}{\textbf{#1}}}
\newcommand{\SpecialStringTok}[1]{\textcolor[rgb]{0.31,0.60,0.02}{#1}}
\newcommand{\StringTok}[1]{\textcolor[rgb]{0.31,0.60,0.02}{#1}}
\newcommand{\VariableTok}[1]{\textcolor[rgb]{0.00,0.00,0.00}{#1}}
\newcommand{\VerbatimStringTok}[1]{\textcolor[rgb]{0.31,0.60,0.02}{#1}}
\newcommand{\WarningTok}[1]{\textcolor[rgb]{0.56,0.35,0.01}{\textbf{\textit{#1}}}}
\usepackage{graphicx}
\makeatletter
\newsavebox\pandoc@box
\newcommand*\pandocbounded[1]{% scales image to fit in text height/width
  \sbox\pandoc@box{#1}%
  \Gscale@div\@tempa{\textheight}{\dimexpr\ht\pandoc@box+\dp\pandoc@box\relax}%
  \Gscale@div\@tempb{\linewidth}{\wd\pandoc@box}%
  \ifdim\@tempb\p@<\@tempa\p@\let\@tempa\@tempb\fi% select the smaller of both
  \ifdim\@tempa\p@<\p@\scalebox{\@tempa}{\usebox\pandoc@box}%
  \else\usebox{\pandoc@box}%
  \fi%
}
% Set default figure placement to htbp
\def\fps@figure{htbp}
\makeatother
\setlength{\emergencystretch}{3em} % prevent overfull lines
\providecommand{\tightlist}{%
  \setlength{\itemsep}{0pt}\setlength{\parskip}{0pt}}
\usepackage{bookmark}
\IfFileExists{xurl.sty}{\usepackage{xurl}}{} % add URL line breaks if available
\urlstyle{same}
\hypersetup{
  pdftitle={Biomarker statistics (Youden{,} NRI{,} decision curves)},
  hidelinks,
  pdfcreator={LaTeX via pandoc}}

\title{Biomarker statistics (Youden, NRI, decision curves)}
\author{}
\date{\vspace{-2.5em}}

\begin{document}
\maketitle

\textbf{Testing labs} use the main template: Hypothesis → Visualise →
Assumptions → Conduct → Conclude.

\subsection{Learning objectives}\label{learning-objectives}

\begin{itemize}
\tightlist
\item
  Compute ROC-AUC, Youden's index, sensitivity, and specificity at the
  optimal cut-point.
\item
  Compute a net reclassification index between two predictive models.
\item
  Construct a decision curve and interpret net benefit at a range of
  threshold probabilities.
\end{itemize}

\subsection{Prerequisites}\label{prerequisites}

Binary classification and ROC curves; logistic regression.

\subsection{Background}\label{background}

A candidate biomarker is not useful until it is tied to a decision.
ROC-AUC summarises discrimination across all thresholds but is
insensitive to \emph{where} on the curve the action happens. Youden's
index (sensitivity + specificity − 1) picks the threshold that maximises
the equal-weighted sum. The net reclassification index (NRI) quantifies
whether a new model reclassifies cases and non-cases in the correct
direction relative to a baseline. Decision curve analysis plots
\emph{net benefit} as a function of the threshold probability, and lets
a reader compare strategies (``treat all'', ``treat none'', ``treat by
model'') across a clinically relevant range.

Discrimination, calibration, and net benefit are three complementary
axes. A biomarker with high AUC that is poorly calibrated can produce
harmful decisions; a perfectly calibrated biomarker with low AUC gives
no useful ranking. Reporting all three keeps the evaluation honest.

Decision curves are not hypothesis tests. They are a principled way to
put a clinical question --- what is the harm-to-benefit ratio of acting
on this prediction? --- into the analysis, and to see how the answer
depends on that ratio.

\subsection{Setup}\label{setup}

\begin{Shaded}
\begin{Highlighting}[]
\FunctionTok{library}\NormalTok{(tidyverse)}
\FunctionTok{library}\NormalTok{(pROC)}
\FunctionTok{library}\NormalTok{(MASS)}
\FunctionTok{set.seed}\NormalTok{(}\DecValTok{42}\NormalTok{)}
\FunctionTok{theme\_set}\NormalTok{(}\FunctionTok{theme\_minimal}\NormalTok{(}\AttributeTok{base\_size =} \DecValTok{12}\NormalTok{))}
\end{Highlighting}
\end{Shaded}

\subsection{1. Hypothesis}\label{hypothesis}

Can a logistic model on \texttt{Pima.tr} (glucose, BMI, age) distinguish
diabetic from non-diabetic patients well enough to support a screening
decision?

\subsection{2. Visualise}\label{visualise}

\begin{Shaded}
\begin{Highlighting}[]
\FunctionTok{ggplot}\NormalTok{(d, }\FunctionTok{aes}\NormalTok{(glu, bmi, }\AttributeTok{colour =}\NormalTok{ type)) }\SpecialCharTok{+} \FunctionTok{geom\_point}\NormalTok{(}\AttributeTok{alpha =} \FloatTok{0.7}\NormalTok{)}
\end{Highlighting}
\end{Shaded}

\subsection{3. Assumptions}\label{assumptions}

Independence of observations; probability of diabetes is monotone in the
linear predictor; no missingness.

\subsection{4. Conduct}\label{conduct}

Fit a simple logistic regression and compute discrimination.

\begin{Shaded}
\begin{Highlighting}[]
\NormalTok{d}\SpecialCharTok{$}\NormalTok{p }\OtherTok{\textless{}{-}} \FunctionTok{predict}\NormalTok{(fit, }\AttributeTok{type =} \StringTok{"response"}\NormalTok{)}
\NormalTok{r  }\OtherTok{\textless{}{-}} \FunctionTok{roc}\NormalTok{(d}\SpecialCharTok{$}\NormalTok{type, d}\SpecialCharTok{$}\NormalTok{p, }\AttributeTok{direction =} \StringTok{"\textless{}"}\NormalTok{, }\AttributeTok{quiet =} \ConstantTok{TRUE}\NormalTok{)}
\FunctionTok{auc}\NormalTok{(r)}
\FunctionTok{coords}\NormalTok{(r, }\StringTok{"best"}\NormalTok{, }\AttributeTok{ret =} \FunctionTok{c}\NormalTok{(}\StringTok{"threshold"}\NormalTok{, }\StringTok{"sensitivity"}\NormalTok{, }\StringTok{"specificity"}\NormalTok{,}
                          \StringTok{"youden"}\NormalTok{), }\AttributeTok{transpose =} \ConstantTok{FALSE}\NormalTok{)}
\end{Highlighting}
\end{Shaded}

NRI against a glucose-only baseline.

\begin{Shaded}
\begin{Highlighting}[]
\NormalTok{p0 }\OtherTok{\textless{}{-}} \FunctionTok{predict}\NormalTok{(fit0, }\AttributeTok{type =} \StringTok{"response"}\NormalTok{)}
\NormalTok{p1 }\OtherTok{\textless{}{-}}\NormalTok{ d}\SpecialCharTok{$}\NormalTok{p}
\CommentTok{\# Continuous NRI}
\NormalTok{case   }\OtherTok{\textless{}{-}}\NormalTok{ d}\SpecialCharTok{$}\NormalTok{type }\SpecialCharTok{==} \StringTok{"Yes"}
\NormalTok{nri\_up }\OtherTok{\textless{}{-}} \FunctionTok{mean}\NormalTok{(p1[case]  }\SpecialCharTok{\textgreater{}}\NormalTok{ p0[case])   }\SpecialCharTok{{-}} \FunctionTok{mean}\NormalTok{(p1[case]  }\SpecialCharTok{\textless{}}\NormalTok{ p0[case])}
\NormalTok{nri\_dn }\OtherTok{\textless{}{-}} \FunctionTok{mean}\NormalTok{(p1[}\SpecialCharTok{!}\NormalTok{case] }\SpecialCharTok{\textless{}}\NormalTok{ p0[}\SpecialCharTok{!}\NormalTok{case])  }\SpecialCharTok{{-}} \FunctionTok{mean}\NormalTok{(p1[}\SpecialCharTok{!}\NormalTok{case] }\SpecialCharTok{\textgreater{}}\NormalTok{ p0[}\SpecialCharTok{!}\NormalTok{case])}
\NormalTok{nri }\OtherTok{\textless{}{-}}\NormalTok{ nri\_up }\SpecialCharTok{+}\NormalTok{ nri\_dn}
\FunctionTok{c}\NormalTok{(}\AttributeTok{nri\_cases =}\NormalTok{ nri\_up, }\AttributeTok{nri\_noncases =}\NormalTok{ nri\_dn, }\AttributeTok{nri\_total =}\NormalTok{ nri)}
\end{Highlighting}
\end{Shaded}

A manual decision curve.

\begin{Shaded}
\begin{Highlighting}[]
\NormalTok{dca }\OtherTok{\textless{}{-}} \FunctionTok{sapply}\NormalTok{(thr, }\ControlFlowTok{function}\NormalTok{(t) \{}
\NormalTok{  treat }\OtherTok{\textless{}{-}}\NormalTok{ p1 }\SpecialCharTok{\textgreater{}}\NormalTok{ t}
\NormalTok{  tp }\OtherTok{\textless{}{-}} \FunctionTok{sum}\NormalTok{(treat }\SpecialCharTok{\&}\NormalTok{  case); fp }\OtherTok{\textless{}{-}} \FunctionTok{sum}\NormalTok{(treat }\SpecialCharTok{\&} \SpecialCharTok{!}\NormalTok{case); N }\OtherTok{\textless{}{-}} \FunctionTok{length}\NormalTok{(case)}
\NormalTok{  tp }\SpecialCharTok{/}\NormalTok{ N }\SpecialCharTok{{-}}\NormalTok{ (fp }\SpecialCharTok{/}\NormalTok{ N) }\SpecialCharTok{*}\NormalTok{ (t }\SpecialCharTok{/}\NormalTok{ (}\DecValTok{1} \SpecialCharTok{{-}}\NormalTok{ t))}
\NormalTok{\})}
\NormalTok{nb\_all }\OtherTok{\textless{}{-}} \FunctionTok{sapply}\NormalTok{(thr, }\ControlFlowTok{function}\NormalTok{(t) \{}
\NormalTok{  tp }\OtherTok{\textless{}{-}} \FunctionTok{sum}\NormalTok{(case); fp }\OtherTok{\textless{}{-}} \FunctionTok{sum}\NormalTok{(}\SpecialCharTok{!}\NormalTok{case); N }\OtherTok{\textless{}{-}} \FunctionTok{length}\NormalTok{(case)}
\NormalTok{  tp }\SpecialCharTok{/}\NormalTok{ N }\SpecialCharTok{{-}}\NormalTok{ (fp }\SpecialCharTok{/}\NormalTok{ N) }\SpecialCharTok{*}\NormalTok{ (t }\SpecialCharTok{/}\NormalTok{ (}\DecValTok{1} \SpecialCharTok{{-}}\NormalTok{ t))}
\NormalTok{\})}
\FunctionTok{tibble}\NormalTok{(}\AttributeTok{threshold =}\NormalTok{ thr, }\AttributeTok{model =}\NormalTok{ dca, }\AttributeTok{treat\_all =}\NormalTok{ nb\_all, }\AttributeTok{treat\_none =} \DecValTok{0}\NormalTok{) }\SpecialCharTok{|\textgreater{}}
  \FunctionTok{pivot\_longer}\NormalTok{(}\SpecialCharTok{{-}}\NormalTok{threshold) }\SpecialCharTok{|\textgreater{}}
  \FunctionTok{ggplot}\NormalTok{(}\FunctionTok{aes}\NormalTok{(threshold, value, }\AttributeTok{colour =}\NormalTok{ name)) }\SpecialCharTok{+} \FunctionTok{geom\_line}\NormalTok{() }\SpecialCharTok{+}
  \FunctionTok{labs}\NormalTok{(}\AttributeTok{x =} \StringTok{"threshold probability"}\NormalTok{, }\AttributeTok{y =} \StringTok{"net benefit"}\NormalTok{)}
\end{Highlighting}
\end{Shaded}

\subsection{5. Concluding statement}\label{concluding-statement}

\begin{quote}
A logistic model using glucose, BMI, and age discriminated diabetic from
non-diabetic patients in \texttt{MASS::Pima.tr} with AUC
\texttt{round(as.numeric(auc(r)),\ 3)}. The Youden-optimal cut-point
occurred at a predicted probability of
\texttt{round(coords(r,\ "best",\ ret\ =\ "threshold",\ transpose\ =\ FALSE){[}1,\ 1{]},\ 2)}.
Adding BMI and age to a glucose-only baseline produced an NRI of
\texttt{round(nri,\ 2)}; the decision curve showed net benefit above
``treat all'' for threshold probabilities from roughly 0.15 to 0.5.
\end{quote}

Decision curves give the clinical context: if the decision to intervene
at, say, p = 0.2 is under discussion, the model is useful; at p = 0.05
or p = 0.7, it is barely distinguishable from treat-all or treat-none.

\subsection{Common pitfalls}\label{common-pitfalls}

\begin{itemize}
\tightlist
\item
  Reporting AUC without calibration or decision curves.
\item
  Computing NRI with a categorical risk cut-point and failing to
  disclose the cut-off.
\item
  Using the same data to develop and evaluate the biomarker (apparent
  performance).
\end{itemize}

\subsection{Further reading}\label{further-reading}

\begin{itemize}
\tightlist
\item
  Pencina MJ, D'Agostino RB Sr, et al.~(2008), \emph{Evaluating the
  added predictive ability of a new marker}.
\item
  Vickers AJ, Elkin EB (2006), \emph{Decision curve analysis}.
\end{itemize}

\subsection{Session info}\label{session-info}

\begin{Shaded}
\begin{Highlighting}[]

\end{Highlighting}
\end{Shaded}

\subsection{See also --- chapter index}\label{see-also-chapter-index}

\begin{itemize}
\tightlist
\item
  \href{../index.Rmd}{Methods in this chapter}
\item
  \href{../../../ch00-find-your-method.Rmd}{Find your method (Chapter
  0)}
\end{itemize}

\end{document}
